\documentclass[10pt,a4paper]{amsart}
\usepackage[margin=19mm]{geometry}
\usepackage{amsmath,amssymb,mathtools}
\usepackage[scr=rsfs,cal=euler]{mathalfa}
\usepackage{xfrac}
\usepackage[unicode,hidelinks,bookmarks=false]{hyperref}
\newcommand{\ie}{i.e.\ }
\newcommand{\cf}{cf.\ }
\newcommand{\resp}{respectively\ }
\newcommand{\Subsection}{\subsection}
\providecommand{\nobreakdash}{\nobreak\hbox{-}}
\setlength{\emergencystretch}{2em}
\multlinegap0pt
\usepackage{dsfont}
\usepackage{amssymb}
\usepackage{mathtools}
\usepackage[shortlabels]{enumitem}
\usepackage{csquotes}
\usepackage{xcolor}
\newtheorem{lem}{Lemma}
\usepackage{cleveref}
\crefname{lem}{Lemma}{Lemmas}
\newcommand{\D}{\mathrm{d}}
\title{Critical Self-Similar Markov Trees}
\author{Nicolas Curien, Xingjian Hu, Dongjian Qian}
\date{Source: arXiv:2603.16408v1 (CC BY 4.0)}
\begin{document}
\maketitle

\noindent\emph{Source and licence.} Critical Self-Similar Markov Trees. Version arXiv:2603.16408v1 (CC BY 4.0), distributed by arXiv under the Creative Commons Attribution 4.0 International licence, http://creativecommons.org/licenses/by/4.0/, as recorded in the arXiv licence metadata for this exact version and shown on the abstract page https://arxiv.org/abs/2603.16408v1. Full licence notice: \texttt{licenses/arxiv-2603.16408-CC-BY-4.0.txt}.\par\smallskip

\noindent\textit{Editorial scope.} Counted unit: the article's lem good (source0.tex lines 1112--1117). The article's complete original proof of it is delivered with it (source0.tex lines 1200--1306, one \texttt{proof} environment of 107 lines, which the article places away from the statement). No mathematics is written, repaired or re-derived: the statement and the proof are the article's own lines, in the article's own notation, and no retained line is substituted or re-flowed.  Retained uncounted as the source-local context the proof needs: lines 399--404; lines 401--403; lines 909--911; lines 1086--1088; lines 1173--1175; lines 1190--1192.  Nothing was omitted from any retained span, so the package carries no omission record.  No retained line contains a citation, so the wrapper carries no bibliography and no bibliography entry.  No retained line contains a figure environment, an image-inclusion command or drawing code, so no image is delivered and the package declares no asset dependency.  Editorial formatting only, in addition: an AMS \texttt{amsart} preamble that restates the article's own declarations and macros used by the retained lines, namely the environments \texttt{lem} on separate counters, together with the standard packages those lines need.  The retained environments print in this document as Lemma 1 in order of appearance, whereas the article prints the same items with the numbering of its own full document.  The complete native source of the article as distributed in the arXiv e-print for this exact version is bundled unchanged as \texttt{sources/196546/source0.tex}.
\begin{lem}\label{lem: many-to-one}
Fix $n \geq 1$, for a generic positive functional $F$ we have
\begin{equation}\label{eq: many-to-one}
\mathbb{E}\bigg[\sum_{|u|=n} (\chi(u))^{\omega_-}F(-\log\chi(u_1), \ldots, -\log\chi(u_n))\bigg] = {\tt E}[F({\hat{\tt S}}_1, \ldots, {\hat{\tt S}}_n)].
\end{equation}
\end{lem}

\begin{equation}
\mathbb{E}\bigg[\sum_{|u|=n} (\chi(u))^{\omega_-}F(-\log\chi(u_1), \ldots, -\log\chi(u_n))\bigg] = {\tt E}[F({\hat{\tt S}}_1, \ldots, {\hat{\tt S}}_n)].
\end{equation}

\begin{equation}\label{eq: rough bound r c}
r_{\gamma}(c)\leq \mathbb{E}\bigg[\int^{\zeta}_0 \mathrm{e}^{\gamma\xi_s}\D s \bigg] = -\frac{1}{\psi(\gamma)}
\end{equation}

\begin{equation}\label{eq: green function r.w.}
(\gamma-\omega_-)\cdot{\tt E}\left[\sum^{{\tt T}^+_b-1}_{n=0} \mathrm{e}^{-(\gamma-\omega_-){\hat{\tt S}}_n}\right]\leq C\mathrm{e}^{(\gamma-\omega_-)b}.
\end{equation}

\begin{equation}\label{eq: frac bound by g}
\frac{(\chi(i))^{\gamma}R_{\gamma}\Big(\frac{c(x)}{\chi(i)}\Big)}{R_{\gamma}(c(x))} \leq C(\chi(i))^{\gamma}\Big(1 + \log_+\Big(\frac{1}{\chi(i)}\Big)\Big) = Cg(\chi(i)).
\end{equation}

\begin{equation}\label{eq: R gamma bounds 2}
(\gamma-\omega_-)R_{\gamma}(c\mathrm{e}^{{\hat{\tt S}}_n})\leq Cc^{\gamma-\omega_-}\mathrm{e}^{(\gamma-\omega_-){\hat{\tt S}}_n}.
\end{equation}

\begin{lem}\label{lem: estimate good}
There exists a constant $C>0$, such that for $(c, A)$ and $\gamma\in(\omega_-, \omega_- + \Delta_0)$, we have
\begin{equation}
\mathbb{E}[(\gamma-\omega_-)^2\tilde{\lambda}^{\gamma}_c(T^c_G)^2]\leq CA^2c^{(2\gamma-\omega_-)}.
\end{equation}
\end{lem}

\begin{proof}[Proof of \cref{lem: estimate good}]
Throughout the proof we allow the universal constant $C>0$ to vary from line to line. We use double indices $u$ and $v$ to interpret $\tilde{\lambda}^{\gamma}_c(T^c_G)^2$ by
\[\tilde{\lambda}^{\gamma}_c(T^c_G)^2 = \bigg(\sum_{u\in\mathbb{G}^c}(\chi(u))^{\gamma}r_{\gamma}\Big(\frac{c}{\chi(u)}\Big)\bigg)\bigg(\sum_{v\in\mathbb{G}^c}(\chi(v))^{\gamma}r_{\gamma}\Big(\frac{c}{\chi(v)}\Big)\bigg). \]

\noindent When expanding the brackets, there will be three cases for $(u, v)$:
\begin{enumerate}
   \item $u = v$;
   \item $u\ne v$, either $u\prec v$ or $v\prec u$ is satisfied;
   \item $u\ne v$, neither $u\prec v$ nor $v\prec u$ is satisfied.
\end{enumerate}
\noindent In formulation, we write
\begin{align*}
\tilde{\lambda}^{\gamma}_c(T^c_G)^2 =& \sum_{u\in\mathbb{G}^c}(\chi(u))^{2\gamma}r_{\gamma}\Big(\frac{c}{\chi(u)}\Big)^2 + 2\sum_{u\in\mathbb{G}^c} \sum_{v\in\mathbb{G}^c} \mathds{1}_{\{u\ne v\}}\mathds{1}_{\{u\prec v\}}(\chi(u))^{\gamma}r_{\gamma}\Big(\frac{c}{\chi(u)}\Big)(\chi(v))^{\gamma}r_{\gamma}\Big(\frac{c}{\chi(v)}\Big)\\
& + \sum_{u\in\mathbb{G}^c} \sum_{v\in\mathbb{G}^c} \mathds{1}_{\{u\ne v\}}\mathds{1}_{\{u\nprec v\}}\mathds{1}_{\{v\nprec u\}} (\chi(u))^{\gamma}r_{\gamma}\Big(\frac{c}{\chi(u)}\Big)(\chi(v))^{\gamma}r_{\gamma}\Big(\frac{c}{\chi(v)}\Big)\\
=& I_1 + I_2 + I_3.
\end{align*}

We first estimate $\mathbb{E}[I_1]$. Using the rough bound $r_{\gamma}(c/\chi(u))\le |\psi(\gamma)|^{-1}$ by \eqref{eq: rough bound r c}, and replacing $\{u\in\mathbb{G}^c\}$ by $\{\chi(u_k)\leq c, k\leq |u|\}$, we see
\[
\mathbb{E}[I_1] \leq |\psi(\gamma)|^{-2}\mathbb{E}\bigg[\sum_{u\in\mathbb{U}}(\chi(u))^{2\gamma}\mathds{1}_{\{\chi(u_k)<c, k\leq |u|\}}\bigg].
\]
By the many-to-one formula \eqref{eq: many-to-one} and then \eqref{eq: green function r.w.}, we have
\begin{equation}\label{eq: proof good I1 final}
\mathbb{E}[I_1] \leq |\psi(\gamma)|^{-2}{\tt E}\bigg[\sum^{{\tt T}_{\log c}^+-1}_{n=0} \mathrm{e}^{-(2\gamma-\omega_-){\hat{\tt S}}_n}\bigg]\leq Cc^{2\gamma-\omega_-}.
\end{equation}

\noindent Notice that $2\gamma-\omega_-$ is bounded away from $0$, so the last inequality is uniform in $\gamma\in(\omega_-, \omega_-+\Delta_0)$.

We continue to estimate $\mathbb{E}[I_2]$. For $u\prec v$ and $u\ne v$, write $v = uw$ for $w\in \mathbb{U}^*$. Rewrite the summation by
\[
\mathbb{E}[I_2] = 2\sum^{\infty}_{n=0}\mathbb{E}\bigg[\sum_{|u|=n}\mathds{1}_{\{u\in\mathbb{G}^c\}} (\chi(u))^{\gamma}r_{\gamma}\Big(\frac{c}{\chi(u)}\Big)\sum_{w\in\mathbb{U}^*} (\chi(uw))^{\gamma}r_{\gamma}\Big(\frac{c}{\chi(uw)}\Big)\mathds{1}_{\{uw\in\mathbb{G}^c\}}\bigg].
\]

\noindent Define the event (the complement of $B_x$)
\begin{equation}\label{eq: def G_x}
    G_x = \bigg\{\sum^{\infty}_{i=1} g\Big(\frac{\chi(i)}{\chi(\varnothing)}\Big)\leq A\Big(\frac{c}{x}\Big)^{\theta} \bigg\}.
\end{equation}

\noindent For each $n\geq 0$, we take conditional expectation with respect to $\mathcal{F}_n$, whence
\begin{equation}\label{eq: proof good I2 1}
\mathbb{E}[I_2] \le 2\sum^{\infty}_{n=0}\mathbb{E}\bigg[\sum_{|u|=n}\mathds{1}_{\{u\in\mathbb{U}^c\}} (\chi(u))^{\gamma}r_{\gamma}\Big(\frac{c}{\chi(u)}\Big)\mathbb{E}_x\Big[\mathds{1}_{G_x}\sum_{w\in\mathbb{U}^*} (\chi(w))^{\gamma}r_{\gamma}\Big(\frac{c}{\chi(w)}\Big)\mathds{1}_{\{w\in\mathbb{G}^{c}\}}\Big]\bigg|_{x=\chi(u)}\bigg].
\end{equation}

\noindent Recall $c(x)=c/x$. Replacing the indicator $\mathds{1}_{\{w\in\mathbb{G}^c\}}$ by $\mathds{1}_{\{w\in\mathbb{U}^c\}}$, by the scaling property, we bound the inner expectation by
\begin{align}  \mathbb{E}_x\bigg[\mathds{1}_{G_x}\sum_{w\in\mathbb{U}^*} (\chi(w))^{\gamma}r_{\gamma}\Big(\frac{c}{\chi(w)}\Big) \mathds{1}_{\{w\in\mathbb{U}^c\}}\bigg] &=x^{\gamma}\mathbb{E}\bigg[\mathds{1}_{G_x}\sum_{w\in\mathbb{U}^*} (\chi(w))^{\gamma}r_{\gamma}\Big(\frac{c(x)}{\chi(w)}\Big)\mathds{1}_{\{w\in\mathbb{U}^{c(x)}\}}\bigg]\nonumber \\
&=x^{\gamma}\mathbb{E}\bigg[\mathds{1}_{G_x}\sum^{\infty}_{i=1} \tilde{\lambda}^{\gamma}_{c(x)}(T^{c(x)}(i))\bigg]. \label{eq: proof good I2 2}
\end{align}

\noindent Take conditional expectation on $\mathcal{F}_1$,  we have
\[\mathbb{E}\bigg[\mathds{1}_{G_x}\sum^{\infty}_{i=1} \tilde{\lambda}^{\gamma}_{c(x)}(T^{c(x)}(i))\bigg] = \mathbb{E}\bigg[\mathds{1}_{G_x}\sum^{\infty}_{i=1} (\chi(i))^{\gamma}R_{\gamma}\Big(\frac{c(x)}{\chi(i)}\Big)\bigg]. \]

\noindent With the inequality \eqref{eq: frac bound by g} and the definition of $G_x$, it follows that
\[\mathbb{E}\bigg[\mathds{1}_{G_x}\sum^{\infty}_{i=1} (\chi(i))^{\gamma}R_{\gamma}\Big(\frac{c(x)}{\chi(i)}\Big)\bigg]\leq CR_{\gamma}(c(x))\mathbb{E}\bigg[\mathds{1}_{G_x}\sum^{\infty}_{i=1} g(\chi(i))\bigg]\leq CAc^{\theta}x^{-\theta}R_{\gamma}(c(x)). \]

\noindent Combining with \eqref{eq: proof good I2 2}, we bound \eqref{eq: proof good I2 1} by
\[\mathbb{E}[I_2]\leq CAc^{\theta}\sum^{\infty}_{n=0}\mathbb{E}\bigg[\sum_{|u|=n}\mathds{1}_{\{u\in\mathbb{G}^c\}} (\chi(u))^{2\gamma-\theta}R_{\gamma}\Big(\frac{c}{\chi(u)}\Big)\bigg].\]

\noindent Applying \cref{lem: many-to-one}, we rewrite
\[
\mathbb{E}[I_2]\leq CAc^{\theta}{\tt E}\bigg[\sum^{{\tt T}^+_{\log c}-1}_{n=0}\mathrm{e}^{-(2\gamma-\theta-\omega_-){\hat{\tt S}}_n} R_{\gamma}(c\mathrm{e}^{{\hat{\tt S}}_n})\bigg].
\]

\noindent The bound of $R_{\gamma}(c)$ in \eqref{eq: R gamma bounds 2} now implies that
\begin{equation}\label{eq: proof good I2 final}
\mathbb{E}[(\gamma-\omega_-)I_2]\leq CAc^{\theta+\gamma-\omega_-} {\tt E}\bigg[\sum^{{\tt T}_{\log c}^+-1}_{n=0}\mathrm{e}^{-(\gamma-\theta){\hat{\tt S}}_n} \bigg]\leq CAc^{2\gamma-\omega_-}.
\end{equation}

\noindent We use \eqref{eq: green function r.w.} in the last inequality. Since $\gamma-\theta$ is bounded away from $0$, the last step is uniform in $\gamma\in(\omega_-, \omega_-+\Delta_0)$.

We finally estimate $\mathbb{E}[I_3]$. Consider the last common ancestor (denoted by $w$) of $u$ and $v$. Write $u = wu'$ and $v = wv'$. Then $u'_1\ne v'_1\in\mathbb{U}^*$. By definition, $w\in\mathbb{G}^c$. We rewrite $\mathbb{E}[I_3]$ as
\[
  \mathbb{E}\bigg[\sum_{w\in\mathbb{G}^c}\sum_{u', v'\in\mathbb{U}^*}\mathds{1}_{\{u'_1\ne v'_1\}}(\chi(wu'))^{\gamma}r_{\gamma}\Big(\frac{c}{\chi(wu')}\Big)(\chi(wv'))^{\gamma}r_{\gamma}\Big(\frac{c}{\chi(wv')}\Big)\mathds{1}_{\{wu'\in\mathbb{G}^c\}}\mathds{1}_{\{wv'\in\mathbb{G}^c\}}\bigg].
\]

\noindent We use the same argument as in the estimate of $\mathbb{E}[I_2]$. For each $n\geq 0$, taking conditional expectation with respect to $\mathcal{F}_n$ and inserting the indicator of $G_{x}$, we see that $\mathbb{E}[I_3]$ is less than
\begin{equation}\label{eq: proof good I3 1}
 \mathbb{E}\bigg[\sum_{w\in\mathbb{U}^c}\mathbb{E}_x\bigg[\mathds{1}_{G_x}\sum_{u, v\in\mathbb{U}^*}\mathds{1}_{\{u_1\ne v_1\}}(\chi(u))^{\gamma}r_{\gamma}\Big(\frac{c}{\chi(u)}\Big)(\chi(v))^{\gamma}r_{\gamma}\Big(\frac{c}{\chi(v)}\Big)\mathds{1}_{\{u\in\mathbb{G}^c\}}\mathds{1}_{\{v\in\mathbb{G}^c\}}\bigg]\bigg|_{x=\chi(w)}\bigg].
\end{equation}

\noindent Replace the indicator $\mathds{1}_{\{u,v\in\mathbb{G}^c\}}$ by $\mathds{1}_{\{u,v\in\mathbb{U}^c\}}$. By the scaling property, we bound the inner expectation by
\begin{align}
&\mathbb{E}_x\bigg[\sum_{u, v\in\mathbb{U}^*}\mathds{1}_{G_x}\mathds{1}_{\{u_1\ne v_1\}}(\chi(u))^{\gamma}r_{\gamma}\Big(\frac{c}{\chi(u)}\Big)(\chi(v))^{\gamma}r_{\gamma}\Big(\frac{c}{\chi(v)}\Big)\mathds{1}_{\{u\in\mathbb{U}^c\}}\mathds{1}_{\{v\in\mathbb{U}^c\}}\bigg]\nonumber\\
&= x^{2\gamma}\mathbb{E}\bigg[\sum_{u, v\in\mathbb{U}^*}\mathds{1}_{G_x}\mathds{1}_{\{u_1\ne v_1\}}(\chi(u))^{\gamma}r_{\gamma}\Big(\frac{c(x)}{\chi(u)}\Big)(\chi(v))^{\gamma}r_{\gamma}\Big(\frac{c(x)}{\chi(v)}\Big)\mathds{1}_{\{u\in\mathbb{U}^{c(x)}\}}\mathds{1}_{\{v\in\mathbb{U}^{c(x)}\}}\bigg]\nonumber\\
&= x^{2\gamma}\mathbb{E}\bigg[\mathds{1}_{G_x}\sum_{i\ne j}\tilde{\lambda}^{\gamma}_{c(x)}(T^{c(x)}_i)\cdot \tilde{\lambda}^{\gamma}_{c(x)}(T^{c(x)}_j)\bigg].\label{eq: proof good I3 2}
\end{align}

\noindent Condition on the first generation, $T^{c(x)}_i$ and $T^{c(x)}_j$ are independent for $i\ne j$. Taking conditional expectation, we see
\begin{align*}
\mathbb{E}\bigg[\mathds{1}_{G_x}\sum_{i\ne j}\tilde{\lambda}^{\gamma}_{c(x)}(T^{c(x)}(i))\cdot \tilde{\lambda}^{\gamma}_{c(x)}(T^{c(x)}(j))\bigg]= \mathbb{E}\bigg[\mathds{1}_{G_x}\sum_{i\ne j}(\chi(i))^{\gamma}R_{\gamma}\Big(\frac{c(x)}{\chi(i)}\Big)(\chi(j))^{\gamma}R_{\gamma}\Big(\frac{c(x)}{\chi(j)}\Big)\bigg].
\end{align*}

\noindent Using the inequality \eqref{eq: frac bound by g} and the definition of $G_x$ \eqref{eq: def G_x} again, the above display is bounded by
\[CR_{\gamma}(c(x))^2\mathbb{E} \bigg[\mathds{1}_{G_x}\sum_{i\ne j} g(\chi(i)) g(\chi(j))\bigg]\leq C A^2 R_{\gamma}(c(x))^2 x^{-2\theta}c^{2\theta}. \]

\noindent Plugging it back to \eqref{eq: proof good I3 2} then back to \eqref{eq: proof good I3 1}, we bound $\mathbb{E}[I_3]$ by
\[\mathbb{E}[I_3]\leq C A^2 c^{2\theta}\sum^{\infty}_{n=0} \mathbb{E}\bigg[\sum_{|w| = n}\mathds{1}_{\{w\in\mathbb{U}^c\}}(\chi(w))^{2\gamma-2\theta}R_{\gamma}\Big(\frac{c}{\chi(w)}\Big)^2\bigg]. \]

\noindent The many-to-one formula (\cref{lem: many-to-one}) implies that
\[\mathbb{E}[I_3]\leq C A^2 c^{2\theta}{\tt E}\bigg[\sum^{{\tt T}_{\log c}^+-1}_{n=0}\mathrm{e}^{-(2\gamma-2\theta-\omega_-){\hat{\tt S}}_n} R_{\gamma}(c\mathrm{e}^{{\hat{\tt S}}_n})^2\bigg].\]

\noindent Using the bounds \eqref{eq: R gamma bounds 2} and \eqref{eq: green function r.w.}, we have
\begin{equation}\label{eq: proof good I3 final}
\mathbb{E}[(\gamma-\omega_-)^2I_3]\leq CA^2c^{2(\theta+\gamma-\omega_-)}{\tt E}\bigg[\sum^{{\tt T}_{\log c}^+-1}_{n=0}\mathrm{e}^{-(\omega_--2\theta){\hat{\tt S}}_n}\bigg]\leq CA^2c^{(2\gamma-\omega_-)}.
\end{equation}

Combining \eqref{eq: proof good I1 final}, \eqref{eq: proof good I2 final} and \eqref{eq: proof good I3 final}, we conclude \cref{lem: estimate good}.
\end{proof}

\end{document}
